% =============================================================================
%  Chapter 11 — Domain Refactor Ontology
% =============================================================================

\chapter{Domain Refactor Ontology}
\label{ch:domain-ontology}

This chapter does not yet present the final implementation of the
\texttt{Domain} refactor.  Its purpose is more foundational: to fix the
\emph{ontology} of the codebase before more invasive changes are made.

The current implementation already contains several useful abstractions:
\texttt{Point}, \texttt{Node}, \texttt{IntegrationPoint},
\texttt{MaterialPoint}, \texttt{NodeSection}, \texttt{MaterialSection}, and
\texttt{SectionGeometry}.  The architectural difficulty is that some of these
types represent \emph{different concepts that happen to coincide
geometrically}, while others represent \emph{the same physical location under
different numerical roles}.  If those roles are not separated explicitly,
\texttt{Domain} becomes the place where geometry, analysis state, quadrature,
and constitutive bookkeeping collapse into a single layer.

The refactor should therefore begin by answering a simple question:

\begin{quote}
What kinds of ``points'' exist in the library, and what exactly does each one
mean?
\end{quote}


% ─────────────────────────────────────────────────────────────────────────
\section{Architectural Diagnosis}
\label{sec:domain-ontology-diagnosis}

At present, the library contains three partially overlapping strata:

\begin{enumerate}[nosep]
  \item \textbf{Geometric entities:} reference cells, mapped element geometry,
        coordinates, frames, section shapes.
  \item \textbf{Discrete unknown carriers:} nodal degrees of freedom,
        constrained DOFs, PETSc indexing.
  \item \textbf{Sampling / constitutive sites:} Gauss points, material points,
        beam sections, shell sections.
\end{enumerate}

The main architectural leak is that \texttt{Node} currently belongs to both
the first and second strata at the same time.  It is a geometric point, but
it is also the owner of nodal DOF indices and PETSc-related mesh indexing.
This is convenient for a first implementation, but it prevents a clean
\texttt{Domain} abstraction because a geometry-only domain should not need to
know anything about nodal unknown layouts.

There is a second leak in the repeated use of the word \emph{section}.  In the
current code, \emph{section} may refer to:

\begin{itemize}[nosep]
  \item a structural cross-section shape (\texttt{SectionGeometry});
  \item a section plane located at a node (\texttt{NodeSection});
  \item a constitutive section sample along a beam or shell
        (\texttt{MaterialSection});
  \item a PETSc \texttt{PetscSection}, which is purely an indexing map and not
        a structural section at all.
\end{itemize}

These are not small naming issues.  They correspond to genuinely distinct
mathematical objects.


% ─────────────────────────────────────────────────────────────────────────
\section{Entity-by-Entity Semantics}
\label{sec:domain-ontology-entities}

\subsection{\texttt{Point<Dim>}}

\texttt{Point<Dim>} is the cleanest entity in the current design.  It is a
coordinate value type: a point in physical space with no identity, no topology,
no DOFs, no material, and no ownership semantics beyond ordinary C++ value
semantics.

\paragraph{What it is.}
A metric object in \(\mathbb{R}^{\mathrm{Dim}}\).

\paragraph{What it is not.}
It is not a mesh vertex, not a quadrature point, not a node in the FEM sense,
and not a constitutive point.

\paragraph{Architectural role.}
\texttt{Point} should remain the lowest-level geometric atom.  Any richer
entity that has coordinates should refine or contain a \texttt{Point}.


\subsection{\texttt{Vertex<Dim>} (proposed and introduced in this step)}

\texttt{Vertex<Dim>} is the missing geometry/topology abstraction.  It is a
\texttt{Point} with a stable mesh identity, but without analysis state.  In
other words:

\[
\texttt{Vertex} = \texttt{Point} + \texttt{mesh identity}.
\]

\paragraph{What it is.}
A vertex of the discrete domain mesh.  It represents a topological 0-cell and
its physical coordinates.

\paragraph{What it is not.}
It is not a carrier of primary unknowns.  It has no nodal DOFs, no constraint
state, and no constitutive variables.

\paragraph{Why it matters.}
If \texttt{Domain} is supposed to live at the geometric level, then it should
own \texttt{Vertex} objects, not \texttt{Node} objects.  The new
\texttt{Vertex.hh} introduced in this iteration fixes the intended target
abstraction without forcing an immediate rewrite of every consumer.


\subsection{\texttt{Node<Dim, Storage>}}

\texttt{Node} is currently defined as
\[
\texttt{Node} = \texttt{Point} + \texttt{id} + \texttt{DoF storage} +
\texttt{PETSc/DMPlex indexing metadata}.
\]

\paragraph{What it is.}
An analysis-layer discrete unknown carrier.  More precisely, it is a nodal site
where one or more primary fields are stored and indexed globally.

\paragraph{What it is not.}
It should not be interpreted as the pure geometric vertex of the domain.

\paragraph{Current status.}
The current implementation uses \texttt{Node} in both senses: as the thing
owned by the domain, and as the thing that carries DOFs.  This is precisely
the coupling that the refactor should remove.

\paragraph{Target interpretation after refactor.}
\[
\texttt{Node} = \texttt{Vertex} + \texttt{field attachment}.
\]
This does not necessarily mean inheritance.  The important point is semantic:
\texttt{Node} belongs to the analysis/model layer, not to the geometry/domain
layer.


\subsection{\texttt{IntegrationPoint<dim>}}

\texttt{IntegrationPoint} is a sampling site used by numerical integration.
It stores physical coordinates, a quadrature weight, and an identifier.

\paragraph{What it is.}
A numerical quadrature sample attached to an element geometry.

\paragraph{What it is not.}
It is not a mesh vertex and not a constitutive model by itself.

\paragraph{Important nuance.}
An integration point represents a \emph{role}, not merely a location.  Two
different entities can occupy the same coordinates and still not be the same
thing.  This is crucial for understanding Lobatto collocation and mixed
formulations.

\paragraph{Potential improvement.}
Eventually, \texttt{IntegrationPoint} should probably be understood as a
\emph{sample site} with:

\begin{itemize}[nosep]
  \item a physical coordinate;
  \item a reference coordinate;
  \item a numerical weight;
  \item possibly the differential measure \(|J|\) or a cached metric.
\end{itemize}

The current implementation already covers the first and third items.


\subsection{\texttt{MaterialPoint<MaterialPolicy>}}

\texttt{MaterialPoint} is the constitutive analogue of a continuum Gauss point.
It binds a \texttt{Material<Policy>} to an \texttt{IntegrationPoint}.

\paragraph{What it is.}
A constitutive state machine attached to a continuum sampling site.

\paragraph{What it is not.}
It is not geometry, not a mesh node, and not a quadrature rule.

\paragraph{Mathematical meaning.}
For a continuum element, the geometric interpolation produces the kinematic
state at a sampling location; the material point then stores and updates the
constitutive history at that location:
\[
\texttt{ElementGeometry} \longrightarrow \texttt{IntegrationPoint}
\longrightarrow \texttt{MaterialPoint}.
\]


\subsection{\texttt{SectionGeometry}}

\texttt{SectionGeometry} is not a point at all.  It is the geometric
description of a structural cross-section: area, moments of inertia, torsional
constant, and shear factors.

\paragraph{What it is.}
Cross-section \emph{shape} data.

\paragraph{What it is not.}
It is not a location, not a node, and not a constitutive point.

\paragraph{Architectural consequence.}
The name is correct, but it must stay clearly separated from both
\texttt{PetscSection} and \texttt{MaterialSection}.


\subsection{\texttt{NodeSection<Dim, Geom>}}

\texttt{NodeSection} combines:

\begin{enumerate}[nosep]
  \item a nodal location;
  \item a local section frame;
  \item a cross-section geometry.
\end{enumerate}

\paragraph{What it is.}
A \emph{section station} located at a node, that is, a nodal section plane with
shape data and orientation.

\paragraph{What it is not.}
It is not a constitutive sampling point unless a formulation explicitly chooses
to collocate section constitutive evaluation at that nodal station.

\paragraph{Why this distinction matters.}
In many beam and shell formulations, the nodal section plane is a geometric
interpolation entity, whereas the constitutive update happens at different
quadrature stations.


\subsection{\texttt{MaterialSection<MaterialPolicy>}}

\texttt{MaterialSection} is the structural analogue of \texttt{MaterialPoint}.
It binds a section constitutive law to an integration point along a beam axis
or shell midsurface.

\paragraph{What it is.}
A constitutive state machine attached to a structural sampling site.

\paragraph{What it is not.}
It is not the same as \texttt{SectionGeometry}, and not the same as
\texttt{NodeSection}.

\paragraph{Mathematical meaning.}
It stores and updates generalized stress-resultants as functions of generalized
strain measures:
\[
\texttt{generalized strains} \mapsto \texttt{section resultants}.
\]

\paragraph{Architectural observation.}
\texttt{MaterialPoint} and \texttt{MaterialSection} are role siblings.  The
first is a continuum constitutive site; the second is a structural
constitutive site.


\subsection{\texttt{PetscSection}}

Although it is not a C++ entity defined in this library, it deserves explicit
mention because it introduces the strongest terminological collision.

\paragraph{What it is.}
A PETSc map from mesh points to numbers of DOFs, offsets, and constraints.

\paragraph{What it is not.}
It is neither a structural section nor a constitutive section.

\paragraph{Architectural rule.}
\texttt{PetscSection} belongs to the solver/indexing infrastructure and should
never be allowed to contaminate the semantics of structural section objects.


% ─────────────────────────────────────────────────────────────────────────
\section{Coincidence Versus Identity}
\label{sec:domain-ontology-coincidence}

The most important conceptual rule for the refactor is the following:

\begin{quote}
Two entities may coincide in space without being the same numerical object.
\end{quote}

This rule resolves the apparent paradoxes raised by Lobatto quadrature and
mixed formulations.

\subsection{Case 1: Gauss--Lobatto beam quadrature}

Suppose a beam formulation uses Gauss--Lobatto integration.  Then the end
quadrature points lie exactly at the beam endpoints.  If the beam also uses
nodal section stations at the endpoints, then:

\begin{itemize}[nosep]
  \item the \emph{nodal section station} at the endpoint;
  \item the \emph{quadrature point} at the endpoint;
  \item the \emph{constitutive section sample} at the endpoint
\end{itemize}

all share the same physical coordinates.

That does \emph{not} mean they should necessarily be represented by the same
object.  Their roles are distinct:

\begin{itemize}[nosep]
  \item the nodal station belongs to interpolation and field attachment;
  \item the quadrature point belongs to integration;
  \item the constitutive section sample belongs to material evolution.
\end{itemize}

The right design principle is therefore:

\begin{quote}
collocation should be represented as a relationship, not as ontological
collapse.
\end{quote}

In practice, this means that a Lobatto endpoint can be modelled as
\emph{three roles sharing one location}, rather than as a single overloaded
object.


\subsection{Case 2: mixed or assumed-strain formulations}

In formulations such as degenerated Timoshenko beams or MITC shell families,
some sampling sites may acquire independent unknowns, tying variables, or
assumed-strain parameters.

Here again, the key is to separate:

\begin{enumerate}[nosep]
  \item the geometric location of the sample;
  \item the fact that it is used in integration;
  \item the fact that additional discrete unknowns live there.
\end{enumerate}

An integration point with independent unknowns is therefore \emph{not}
``a node in disguise''; rather, it is a sampling site that additionally plays
the role of a field carrier.

\paragraph{Recommended interpretation.}
The library should treat this as:
\[
\texttt{SampleSite} + \texttt{FieldAttachment},
\]
not as a reason to redefine \texttt{IntegrationPoint} itself as a node.

This preserves both conceptual clarity and extensibility.  The same sampling
site can later carry:

\begin{itemize}[nosep]
  \item a constitutive law only;
  \item constitutive law plus internal variables;
  \item constitutive law plus independent element-local DOFs;
  \item tying or assumed-strain data without constitutive history.
\end{itemize}


% ─────────────────────────────────────────────────────────────────────────
\section{Proposed Ontology for the Refactor}
\label{sec:domain-ontology-proposal}

The cleanest layered ontology for \texttt{fall\_n} is the following.

\subsection{Layer 0: pure geometry/topology}

\begin{itemize}[nosep]
  \item \texttt{Point<Dim>}:
        coordinates only.
  \item \texttt{Vertex<Dim>}:
        point + mesh identity.
  \item \texttt{ElementGeometry<Dim>}:
        mapped geometry, reference interpolation, quadrature mapping, face
        extraction.
  \item \texttt{Domain<Dim>}:
        owns vertices, element geometries, boundary groups, and PETSc mesh
        materialisation support, but no DOFs and no constitutive state.
\end{itemize}

\subsection{Layer 1: discrete field attachment}

\begin{itemize}[nosep]
  \item \texttt{Node}:
        a field-bearing attachment to a vertex, carrying nodal DOFs and
        constraints.
  \item future internal field sites:
        element-local unknown carriers for mixed methods or assumed-strain
        formulations.
  \item \texttt{PetscSection}:
        global indexing metadata for those field sites.
\end{itemize}

\subsection{Layer 2: sampling and constitutive sites}

\begin{itemize}[nosep]
  \item \texttt{IntegrationPoint}:
        numerical sample site.
  \item \texttt{MaterialPoint}:
        continuum constitutive site.
  \item \texttt{MaterialSection}:
        structural constitutive site.
  \item future tying or assumed-strain sites:
        sampled kinematic fields that may or may not have constitutive state.
\end{itemize}

\subsection{Layer 3: structural section semantics}

\begin{itemize}[nosep]
  \item \texttt{SectionGeometry}:
        cross-section shape properties.
  \item \texttt{NodeSection}:
        nodal section station = node + frame + section geometry.
  \item section constitutive relations
        (\texttt{TimoshenkoBeamSection3D}, \texttt{MindlinShellSection},
        \texttt{FiberSection}):
        map generalized strains to resultants.
\end{itemize}

This layering produces a clean separation between:

\begin{enumerate}[nosep]
  \item \emph{where something is};
  \item \emph{how it is sampled};
  \item \emph{what unknowns live there};
  \item \emph{what constitutive law is evaluated there}.
\end{enumerate}


% ─────────────────────────────────────────────────────────────────────────
\section{Implications for \texttt{Domain}}
\label{sec:domain-ontology-domain}

Under this ontology, \texttt{Domain} should evolve toward the following
responsibilities:

\begin{enumerate}[nosep]
  \item ownership of mesh vertices;
  \item ownership of element geometry objects;
  \item storage of boundary groups and physical labels;
  \item translation to \texttt{DMPlex};
  \item topology and geometry queries.
\end{enumerate}

And it should explicitly \emph{not} own:

\begin{enumerate}[nosep]
  \item nodal DOF layouts;
  \item constitutive points;
  \item material state;
  \item structural section stations;
  \item solver vectors or matrices.
\end{enumerate}

This has a direct consequence for the current code:

\begin{quote}
\texttt{Domain} owning \texttt{Node} is architecturally temporary, not final.
\end{quote}

The first code step taken in this iteration is therefore the introduction of
\texttt{Vertex<Dim>} as the future geometry-level entity.


% ─────────────────────────────────────────────────────────────────────────
\section{HPC and Compile-Time Guidance}
\label{sec:domain-ontology-hpc}

The user requirement for the refactor is not only architectural cleanliness but
also efficiency in a high-performance context.  That requirement strongly
favours a few design rules.

\subsection{Prefer static structure in hot paths}

The current library already uses static polymorphism effectively in the element
kernel.  That should continue.  In particular:

\begin{itemize}[nosep]
  \item shape functions, quadrature rules, and reference-cell topology should
        remain compile-time objects whenever possible;
  \item constitutive updates should stay inside element-local loops with no new
        virtual dispatch;
  \item small fixed-size Eigen matrices are still the right choice for local
        operators.
\end{itemize}

\subsection{Keep \texttt{Domain} simple, not hyper-generic}

The right way to favour HPC here is \emph{not} to template \texttt{Domain} on
every imaginable backend or policy.  That would increase compile times,
fragment code paths, and create speculative abstraction.

The efficient choice is:

\begin{itemize}[nosep]
  \item keep \texttt{Domain} concrete and PETSc-oriented;
  \item separate concepts cleanly;
  \item minimise includes and cross-layer dependencies;
  \item keep backend-specific code concentrated in a few classes.
\end{itemize}

\subsection{Avoid type collapse when roles overlap}

Role overlap should be represented through composition and index maps rather
than inheritance-heavy unions of responsibilities.

This means:

\begin{itemize}[nosep]
  \item a Lobatto endpoint may be both nodal station and quadrature sample, but
        that is best represented by shared coordinates or shared IDs, not by
        redefining one type to become the other;
  \item a mixed formulation with internal DOFs at sample sites should introduce
        a dedicated field attachment for those sites, not redefine every
        integration point in the library as a node.
\end{itemize}

\subsection{Recommended storage philosophy}

For the next stages of the refactor, the following storage strategy is the most
coherent with the current codebase:

\begin{itemize}[nosep]
  \item \texttt{Vertex}: compact aggregate, trivially copyable if possible;
  \item \texttt{Node}: inline or SBO DOF storage as already implemented;
  \item \texttt{IntegrationPoint}: compact POD-like sample record;
  \item constitutive points: element-local vectors, because they are naturally
        owned by elements and evolved elementwise;
  \item PETSc indexing: stored in model/solver-side data structures, not in the
        geometric domain.
\end{itemize}

This preserves cache locality while keeping compile-time structure strong and
runtime dispatch localized.


% ─────────────────────────────────────────────────────────────────────────
\section{Migration Plan}
\label{sec:domain-ontology-migration}

The refactor should proceed incrementally.

\subsection*{Phase 1: fix ontology}

\begin{itemize}[nosep]
  \item introduce \texttt{Vertex};
  \item document the role separation;
  \item stop using ``section'' ambiguously in new code.
\end{itemize}

\subsection*{Phase 2: decouple geometry from nodal unknowns}

\begin{itemize}[nosep]
  \item move \texttt{Domain} toward \texttt{std::vector<Vertex<dim>>};
  \item make \texttt{ElementGeometry} depend on a geometry-only vertex concept
        rather than directly on \texttt{Node};
  \item relocate DOF assignment and constraints fully to \texttt{Model}.
\end{itemize}

\subsection*{Phase 3: make field attachments explicit}

\begin{itemize}[nosep]
  \item define how primary nodal unknowns are attached to vertices;
  \item introduce dedicated internal field sites for mixed methods when needed;
  \item keep collocation and coincidence as relationships, not type collapse.
\end{itemize}

\subsection*{Phase 4: unify constitutive-site patterns}

\begin{itemize}[nosep]
  \item consider a shared abstraction behind \texttt{MaterialPoint} and
        \texttt{MaterialSection};
  \item preserve policy-based constitutive dispatch and element-local storage.
\end{itemize}


% ─────────────────────────────────────────────────────────────────────────
\section{Conclusion}
\label{sec:domain-ontology-conclusion}

The central conclusion of this chapter is simple:

\begin{quote}
The library does not suffer from having too many abstractions; it suffers from
having a few abstractions that currently play more than one role.
\end{quote}

The remedy is therefore not more runtime polymorphism, but better ontological
separation:

\begin{itemize}[nosep]
  \item \texttt{Point} for coordinates;
  \item \texttt{Vertex} for mesh geometry/topology;
  \item \texttt{Node} for field attachment and DOFs;
  \item \texttt{IntegrationPoint} for numerical sampling;
  \item \texttt{MaterialPoint}/\texttt{MaterialSection} for constitutive state;
  \item \texttt{SectionGeometry}/\texttt{NodeSection} for structural section
        geometry and nodal section stations.
\end{itemize}

Once this distinction is enforced, the \texttt{Domain} refactor becomes much
less ambiguous.  The domain can stay geometrically clean, the model can own the
discrete unknowns, and the element layer can continue to exploit strong
compile-time structure without sacrificing conceptual clarity.
